% from ega3-7-fr.tex
\oldpage[III]{196}
\begin{enumerate}
  \item[(ii)] Supposons en outre que $A$ soit local et réduit, et soit $k$ son
    corps résiduel. Alors, pour que $T_p$ soit exact, il faut et il suffit que
    l'on ait
    \[
      \operatorname{rang}_k T_p(k)
      =
      \operatorname{rang}_{k_i}T_p(k_i)
      \qquad\text{pour }1\leq i\leq r.
      \tag{7.6.10.3}\label{III.7.6.10.3.fr}
    \]
\end{enumerate}
\end{corollary}

\begin{proof}
(i) est immédiat puisque tout voisinage de $x$ contient un des $\mathfrak p_i$,
et il suffit d'appliquer la définition de la semi-continuité. D'autre part, si
$A$ est local, le seul voisinage dans $\operatorname{Spec}(A)$ de l'idéal
maximal $\mathfrak m$ est $\operatorname{Spec}(A)$ tout entier, donc on a
\[
  d_p(x)\leq\operatorname{rang}_k T_p(k)
\]
pour tout $x\in\operatorname{Spec}(A)$; cette relation, jointe à (i), montre que
la condition (7.6.10.3) entraîne que $d_p(x)$ est constante dans
$\operatorname{Spec}(A)$, et par suite que $T_p$ est exact en vertu de
(7.6.9, (ii)); la réciproque est évidente en vertu de (7.6.9, (ii)).
\end{proof}

\begin{remark}[7.6.11]
On peut se demander si l'assertion de (7.6.9, (i)) ne peut être renforcée par
l'inégalité
\[
  \operatorname{rang}_{\kappa(x)}T_p(\kappa(x))
  \geq
  \operatorname{rang}_{\kappa(x)}
    \bigl(T_p(A)\otimes_A\kappa(x)\bigr)
  \tag{7.6.11.1}\label{III.7.6.11.1.fr}
\]
pour tout $x\in\operatorname{Spec}(A)$, qui a lieu effectivement lorsque $A$
est un anneau de valuation discrète et $x$ son idéal maximal (7.6.5).
Bornons-nous au cas où $A$ est un anneau local noethérien, d'idéal maximal
$\mathfrak m$ et de corps résiduel $k$. Alors, les conditions suivantes sont
équivalentes :
\begin{enumerate}
  \item[a)] Pour tout complexe $P_\bullet$ de $A$-modules plats de type fini,
    on a
    \[
      \operatorname{rang}_k T_p(k)
      \geq
      \operatorname{rang}_k\bigl(T_p(A)\otimes_A k\bigr)
      \quad\text{pour tout entier }p.
      \tag{7.6.11.2}\label{III.7.6.11.2.fr}
    \]
  \item[b)] Pour tout $A$-module $M$ de type fini, on a
    \[
      \operatorname{rang}_k(M\otimes_A k)
      \geq
      \operatorname{rang}_k(\check M\otimes_A k).
      \tag{7.6.11.3}\label{III.7.6.11.3.fr}
    \]
  \item[c)] Pour tout $A$-module $N$ de type fini, on a
    \[
      \operatorname{rang}_k\bigl(\operatorname{Tor}_1^A(N,k)\bigr)
      \geq
      \operatorname{rang}_k\bigl(\operatorname{Tor}_2^A(N,k)\bigr).
      \tag{7.6.11.4}\label{III.7.6.11.4.fr}
    \]
\end{enumerate}
On notera qu'il revient au même, par décalage (M, V, 7.2), de dire que l'on a,
pour tout $i\geq1$,
\[
  \operatorname{rang}_k\bigl(\operatorname{Tor}_i^A(N,k)\bigr)
  \geq
  \operatorname{rang}_k\bigl(\operatorname{Tor}_{i+1}^A(N,k)\bigr).
  \tag{7.6.11.5}\label{III.7.6.11.5.fr}
\]
Donnons rapidement quelques indications sur la démonstration. Pour voir que a)
entraîne b), on considère une suite exacte
\[
  L_1\xrightarrow{d}L_0\longrightarrow M\longrightarrow0
\]
où $L_0$ et $L_1$ sont libres de type fini, et on applique a) au complexe
\[
  P_1\xrightarrow{{}^t d}P_0,
  \qquad
  P_0=\check L_1,
  \quad
  P_1=\check L_0,
\]
les autres termes étant nuls; on a alors $T_1(A)=\check M$ et
\[
  T_1(k)
  =
  \operatorname{Hom}_A(M,k)
  =
  \operatorname{Hom}_k(M/\mathfrak mM,k),
\]
autrement dit $T_1(k)$ est le dual de l'espace vectoriel $M\otimes_A k$, et a
donc même rang que ce dernier. Pour prouver que b) entraîne c), nous établissons
d'abord le lemme suivant :

\begin{lemma}[7.6.11.6]
Étant donné un complexe
\[
  \cdots\longrightarrow0\longrightarrow
  P_1\xrightarrow{d}P_0
  \longrightarrow0\longrightarrow\cdots
\]
de $A$-modules plats, on a une suite exacte
\[
  0\longrightarrow\operatorname{Tor}_2^A(Z'_0,k)
  \longrightarrow T_1(A)\otimes_A k
  \longrightarrow T_1(k)
  \longrightarrow\operatorname{Tor}_1^A(Z'_0,k)
  \longrightarrow0.
  \tag{7.6.11.7}\label{III.7.6.11.7.fr}
\]
\end{lemma}

\begin{proof}
